\documentclass[10pt]{article}

\usepackage[utf8]{inputenc}

\usepackage[T1]{fontenc}

\usepackage{amsmath}

\usepackage{amsfonts}

\usepackage{amssymb}

\usepackage[version=4]{mhchem}

\usepackage{stmaryrd}

\usepackage{bbold}

\usepackage{graphicx}

\usepackage[export]{adjustbox}

\graphicspath{ {./images/} }

\usepackage{mathrsfs}



\begin{document}

инвариантных торов при изменении параметров динамической системы. Рождение (или смерть при обратном изменении параметров) гладкого тора, неустранимым образом встречающееся в конечнопараметрич. семействах, возможно лишь при смене устойчивости положения равновесия или предельного цикла. Негладкие торы могут рождаться и умирать и в однопараметрич. семействах в результате нелокальных бифуркаций. Б. же смены устойчивости торов (иногда сопровождаемые возникновением торов более высокой размерности) возможны лишь в бесконечно вырожденных семействах динамич. систем



Jum.: [1] Che nc iner A., I ooss G., "Arch. Ration. Mech. and Anal.», 1979, v. 69, № 2, p. 109-98. БиФУРКАция р ожден и я тора - перестройка в фазовом пространстве одно- или многопараметрического семейства динамических систем, в результате которой возникает и устойчиво существует в некоторой области параметров инвариантный тор.



Пусть семейство динамич. систем непрерывно зависит от параметра $\varepsilon$, принадлежащего топологич. пространству $E$; $G \subset E-$ область существования тора; $\varepsilon_{0} \in \bar{G}-$ бифуркационное значение параметра, отвечающее его рождению; $L=\underset{\varepsilon \rightarrow \varepsilon_{0}}{\overline{\mathrm{lt}}} \mathbb{\pi}_{\varepsilon}-$ верхний топологич. предел (множество всех предельных точек) семейства торов $\left\{T_{\varepsilon}\right\}$. Если $L-$ положение равновесия или предельный цикл, то тор рождается в результате локальной Б., в противном случае - нелокальной. Рождение тора из положения равновесия неустранимым образом встречается в двухпараметрич. семействах динамич. систем (см. таблицу).



\begin{center}

\includegraphics[max width=\textwidth]{2023_07_08_b28afe8d0df10418339bg-1}

\end{center}



В первом столбце таблицы указана размерность центрального многообразия $W^{c}$ положения равновесия $L$, претерпевающего бифуркацию, во втором - приведены семейства динамич. систем на соответствующих центральных многообразиях ( $r_{1,2}$ - амплитудные переменные, циклич. переменные опущены, а $f$ - однородный многочлен 2-й степени от $\left.r_{1}^{2}, r_{2}^{2}, \varepsilon_{1}\right)$, содержащие основную информацию о структуре и Б. динамич. систем в окрестности $L$ (так наз. главные семейства).



Рождение тора при смене устойчивости предельного цикла встречается неустранимым образом уже в однопараметрич. семействах динамич. систем. Для типичного семейства мультипликатор $L$ при $\varepsilon=0$ равен $e^{i \omega}, \omega \neq 2 \pi p / q ; p, q-$ целые, $q \in\{1 ; 2 ; 3 ; 4\}$, а отображение Пуанкаре на центральном многообразии в нек-рых полярных координатах имеет вид:



\[

(r, \theta) \mapsto\left(r+a r^{3}+o\left(r^{3}\right), \theta+\omega+b r^{2}+o\left(r^{2}\right)\right), \quad a \neq 0 .

\]



При изменении є рождается двумерный тор - устойчивый при $a<0$ и неустойчивый при $a>0$.



Инвариантные, вообще говоря, негладкие двумерные торы могут возникать также и в результате нелокальных бифуркаций - напр., вырождаясь из странного аттрактора (при обратном изменении параметра негладкий тор разрушается, порождая странный аттрактор) или склеиваясь из сепаратрис появляющихся предельных циклов седло-узлового типа. При изменении параметров эти торы могут приобретать и увеличивать гладкость и становиться эргодическими.



Лum.: [1] Марсден Дж., Мак Кракен М., Бифуркация рождения цикла и ее приложения, пер. с англ., М., 1980; [2] Арнольд В.И., АфраймовичВ.С., Ильяшенко Ю.С., Ш ильников Л. П., Теория бифуркаций, в кн.: Итоги науки и техники. Современные проблемы математики. Фундаментальные направления, М., 1986, т. 5, с. 3-218.



БИХАРАКТЕРИСТИКА - см. Разностное уравнение; асимптотика решения, Характеристика дифференциального оператора.



БКШ-МОДЕЛЬ, Б а р д и на - Ку п е ра - Ш р и ф ф е р а модель, - модель для микроскопического объяснения сверхпроводимости (см. [1]). Исходит из Фрелиха модели электрон-фононного взаимодействия. Взаимодействие $H_{I}$ ведет к эффективному притяжению электронов при разности энергий между электронными состояниями



\[

\left|\varepsilon_{\vec{k}}-\varepsilon_{\vec{k}+q}\right|<\hbar \omega_{q}

\]



Критерий существования сверхпроводящей фазы - доминирование $H_{I}$ над экранированным кулоновским отталкиванием.



Л. Купер [2] показал, что основное состояние системы из двух электронов с полным нулевым импульсом, взаимодействующих через постоянные отрицательные матричные элементы в малой ячейке над поверхностью Ферми, отделено от континуума энергетич. щелью. Это позволило для получения основного состояния БКШ-м. вариационным способом. использовать редуцированную задачу с включением лишь конфигураций, занятых парами электронов с противоположными спинами и импульсами:



\[

\begin{gathered}

H_{\text {БКШ }}=\sum_{\vec{k}, \sigma^{\prime}}\left(\varepsilon_{\vec{k}}-\mu\right) a_{\vec{k}, \sigma}^{+} a_{\vec{k}, \sigma}- \\

-(2 V)^{-1} \sum_{\vec{k}, \overrightarrow{k^{\prime}}, \sigma, \sigma^{\prime}} \mathscr{V}\left(k, k^{\prime}\right) a_{\vec{k}, \sigma}^{+} a^{+} \vec{k},-\sigma-\overrightarrow{k^{\prime},-\sigma^{\prime}} a_{k^{\prime}, \sigma^{\prime}}

\end{gathered}

\]



В БКШ-М. взаимодействие $\gamma\left(k, k^{\prime}\right)$ факторизуется $\vartheta\left(k, k^{\prime}\right)=$ $=\lambda(k) \lambda\left(k^{\prime}\right) g, g>0$, и отлично от нуля лишь в окрестности поверхности Ферми:



\[

\lambda(k)= \begin{cases}\operatorname{sign} \sigma, & \left|\varepsilon_{\vec{k}}-\mu\right| \leqslant \hbar \omega_{D}, \\ 0, & \left|\varepsilon_{\vec{k}}-\mu\right|>\hbar \omega_{D} .\end{cases}

\]



Основное состояние сверхпроводящего типа не возникает в теории возмущений ни в каком порядке. Его энергия выражается через плотность состояний на поверхности Ферми $N(0)$ :



\[

W_{0}=-2 N(0)\left(\hbar \omega_{D}\right)^{2} /[\exp (2 / N(0) g)-1],

\]



и в пределе слабой связи вычислена критич. температура



\[

\theta_{c}=\hbar \omega_{D} \exp (-1 / N(0) g) \text {. }

\]



В случае сильной связи существование малого параметра в теории обеспечивается теоремой Мигдала. Спектр элементарных возбуждений над основным состоянием имеет вид:



\[

E_{k}=\left(\varepsilon_{k}^{2}+\varepsilon_{0}^{2}\right)^{1 / 2}

\]



где



\[

\varepsilon_{0}=\hbar \omega_{D}(\sin h[1 / N(0) g])^{-1}

\]



\begin{itemize}

  \item щель при нуле температур.

\end{itemize}



\end{document}